\subsection[The pair (1,3): a sixth moment]{The pair $(1,3)$: a sixth moment}\label{loc:S9}
Let $P=qf_3$, $R=r+s$ and $m=1-d/2$. Exact recursion gives
\[
 P\le(m+R/2)^2,\quad\E[pR]=-d-\delta_{2},\quad
 \E R^2=\delta_{2}+\delta_{3},\quad\|R\|_\infty^2\le d+\delta_{2}.
\]

\proofpart{The large-variance range}
Let $c=pq+2r$ be the nonnegative cycle marginal and
$\mathcal T(x)=\E_{y,z}[C(x,y,z)(g(y)+g(z))/2]$.
Symmetrizing gives $\mathcal T=K(g^2)+s-gr$ and $|\mathcal T|\le c$.
The exact path decomposition is
\[
 f_3=p^3/6+(p+2)c/6+(Ac+\mathcal T)/3.
\]
Its remainder after $p^3/6$ is nonnegative, since it is at least
$pc/6+Ac/3$. Put $a_0=qp^3/6\le9/32$.
Using $\mathcal T\le c$ and $T(pq)=q-d+K(g^2)\le q$ in the adjoint integral gives
\[
 \E[p(P-a_0)]\le(113/96)(1-3d).
\]
Here the scalar coefficient is $(7-g-5g^2-g^3)/6\le113/96$.
The row bound gives $P<7/5$ globally. The chord of $x^6$ on
$[a_0,7/5]$ has slope at most
$(7/5)^6/(7/5-9/32)<27/4$.
Thus $\E[pP^6]<1/1000+(3051/384)(1-3d)<1$ for $d\ge7/24$.

\proofpart{Centered variance and a Hermite upper bound}
Now $0<d\le7/24$. Write $a=\delta_{2}/d\in[d,\theta]$, $c_0=45/32$.
Centering $R^2$ gives
\[
 \E[pR^2]\le B_a=c_0ad+d^{3/2}\sqrt{c_0a(1-\theta a)}.
\]
The variance substitution is valid because $c_0a/(1+a)\le\theta<1/2$;
also $B_\theta-B_a\ge c_0d(\theta-a)$.
Set
\begin{gather*}
 t=\sqrt{c_0d},\quad k=16/45,\quad b=13/32,\quad c=5/12,\quad
 v=b+ct,\quad D=1+2t+v,\\
 s_0=-t(t+v)/(1+t),\quad w=(v-t^2)/D,\quad
 U=1+t/2-kt^2,\quad W=1-kt^2+s_0/2.
\end{gather*}
The actual $t$ range is contained in $[0,13/20]$.
Here $v<1$, $v-t^2>0$, $-9/16\le s_0\le0$,
$m-t/2\ge61/120>0$ and $m+s_0/2\ge53/96$.
The quadratic Hermite interpolant of $\phi(x)=(m+x/2)^{12}$ at
the double node $s_0$ and upper node $t$ majorizes $\phi$ on $[-t,t]$,
because $\phi'''\ge0$ and the remainder has sign
$(x-s_0)^2(x-t)\le0$. Write it as $\alpha x^2+\beta x+\gamma$.
Then $\alpha\ge\phi''(s_0)/2>0$ and
\[
 \beta/\alpha\le4m/11-20s_0/11\le61/44<c_0.
\]
Hence the combined mean/variance substitution, including a possibly
negative $\beta$, has the correct sign:
\[
 \alpha B_a-\beta(1+a)d+\gamma
 \le\alpha t^2v-\beta t^2+\gamma.
\]
The two-node law on $t,s_0$ of masses $w,1-w$ has mean $-t^2$ and second
moment $t^2v$. Consequently
$\E[pP^6]\le H(t)=wU^{12}+(1-w)W^{12}$.

\proofpart{The full small interval}
For $t\le1/3$, another Hermite interpolant, at $-t/2$ (double) and $t$,
gives $H\le Q$ with
\[
 Q=\omega\phi(t)+(1-\omega)\phi(-t/2)+h\phi'(-t/2),\quad
 \omega=7/24-7t/27,\quad h=t/16-11t^2/18.
\]
For $x^j$ its expectation is $t^j(A_j+B_jt)$, where
\[
 A_j=7/24+(17-3j)(-1/2)^j/24,\qquad
 B_j=-7/27+(7+33j)(-1/2)^j/27.
\]
For $j\ge5$, $A_j\le3/10$, $B_j\le0$.
For even $j\ge6$ the latter follows from $7+33j\le7\,2^j$;
odd $j$ are immediate. Let $x=kt^2$, and let $E_{n,r}(x)$ be the
degree-$r$ Taylor polynomial of $(1-x)^n$ at zero.
Even truncations are upper bounds and odd truncations lower bounds.
The following finite table specifies all retained terms, including their
required truncation direction:
\begin{center}\small
\begin{tabular}{@{}cc@{}}\toprule
term in $Q$&bound for its factor $m^n$\\\midrule
$m^{12}$&$E_{12,4}$\\
$-6t^2m^{11}$&$E_{11,3}$ (lower)\\
$(33/2)t^2(b+ct)m^{10}$&$E_{10,2}$\\
$(55/8)t^3m^9-(165/8)t^4m^9$&$E_{9,2}$; $E_{9,1}$ (lower)\\
$\widehat A_4t^4m^8+\widehat B_4t^5m^8$&$E_{8,2}$\\
$\widehat A_5t^5m^7-\widehat N_5t^6m^7$&$E_{7,2}$; $E_{7,1}$ (lower)\\
$\widehat A_6t^6m^6-\widehat N_6t^7m^6$&$1$; $E_{6,1}$ (lower)\\
$j=7$&discard $B_7\le0$, then $m^5\le1$\\
$j\ge8$&$A_j\le3/10$, $B_j\le0$, $m^{12-j}\le1$\\\bottomrule
\end{tabular}
\end{center}
The coefficient symbols here are scalar constants:
\begin{gather*}
 \widehat A_4=19305/2048,\ \widehat B_4=495/256,\ \widehat A_5=3663/512,\ \widehat N_5=363/32,\\
 \widehat A_6=34419/8192,\ \widehat N_6=2079/1024,\ \widehat A_7=7425/4096.
\end{gather*}
Substitution gives coefficient upper bounds for powers $t^2,\ldots,t^9$
equal, respectively, to
\[
 -57/16,\quad55/4,\quad-3,\quad-37,\quad19,\quad50,\quad116,\quad32.
\]
For example the first exact coefficient is $-3421/960<-57/16$;
the other entries follow from the finite table and the displayed constants.
The entire $j=10,11,12$ positive tail is at most
$2449t^{10}/122880<t^{10}/50$.
Thus
\begin{align*}
 Q&\le1+t^2[-57/16+55t/4-3t^2-37t^3
                          +19t^4+50t^5+116t^6+32t^7+t^8/50]\\
 &\le1+t^2[-57/16+55t/4-3t^2-20t^3]
 \le1-23t^2/432.
\end{align*}
The second line uses $t\le1/3$ to bound the last five positive terms by
$17t^3$; the final bracket increases to $-23/432$ at $1/3$.

\proofpart{A rational maximum bracket on the remaining interval}
The function $U$ increases and $w$ decreases on $[0,13/20]$.
For $A=wU^{12}$,
\[
 A'=U^{11}N(t)/D^2,\quad
 N=4657/1536+73t/48-84049t^2/4320-2441t^3/216+580t^4/27.
\]
The exact integer polynomial $69120N$ is
$209565+105120t-1344784t^2-781120t^3+1484800t^4$.
It is positive up to $1/10$; thereafter
$N'<73/48-(4623/200)t<0$.
At $17/40$, $-1/200<N<0$; on $[53/125,17/40]$ we have $N'<-7$,
so $N(53/125)>1/500$.
Thus the unique maximum of $A$ lies in that rational bracket.
Its value satisfies
\[
 A\le w(53/125)U(17/40)^{12}
 <(83/500)(263/50)=21829/25000<7/8.
\]
The certificates are $w(53/125)<83/500$ (the reduced numerator is
$-783/4000$), $U(17/40)=20669/18000$ and
$(20669/18000)^3<53/35$, $(53/35)^4<263/50$.
For $B=(1-w)W^{12}$, $W'\le-b/2$ and $W\le1$ give
$B'/B\le2-6b=-7/16<0$.
At $1/3$, $w=125/637$, $W=88199/103680<23/27$,
and $(23/27)^3<5/8$, whence $B(1/3)<625/5096<1/8$.
Therefore $H<1$ on $[1/3,13/20]$.
The two scalar intervals cover $0<d\le7/24$; the cycle argument covers
$d\ge7/24$. Hence $\E[p(qf_3)^6]\le1$ in the full finite model.
